% !TeX root = ../../main.tex
\subsection{Thiết kế mức logic}

\subsubsection{Lựa chọn lược đồ}
\CacBuoc
\begin{enumerate}
\item So sánh star schema, snowflake schema và fact constellation.
\item Đánh giá từng lược đồ theo độ dư thừa, độ phức tạp truy vấn và số bảng sự kiện cần có.
\item Chọn lược đồ và nêu lý do.
\end{enumerate}
\TaiLieu{Ch3 tr.31--37.}
\KetQua{Bảng so sánh và lược đồ được chọn kèm lý do.}
\PhanNhom{lược đồ được chọn}

\subsubsection{Thiết kế bảng sự kiện}
\CacBuoc
\begin{enumerate}
\item Đặt tên từng bảng sự kiện và ghi rõ độ mịn.
\item Liệt kê các khóa ngoại đến bảng chiều.
\item Liệt kê các độ đo, kiểu dữ liệu và nguồn dữ liệu.
\item Xác định khóa chính và cách lưu các mã nghiệp vụ như \code{order\_id}.
\end{enumerate}
\TaiLieu{Ch3 tr.24, 31--32; Ch2 tr.11.}
\KetQua{Bảng mô tả cột của từng bảng sự kiện.}
\PhanNhom{mô tả các bảng sự kiện}

\subsubsection{Thiết kế bảng chiều}
\CacBuoc
\begin{enumerate}
\item Thiết kế từng bảng chiều: tên cột, kiểu dữ liệu, nguồn và ý nghĩa.
\item Định nghĩa khóa thay thế (surrogate key) và giữ lại khóa nghiệp vụ.
\item Thiết kế chiều thời gian theo các phân cấp cần phân tích.
\item Chọn cách xử lý thay đổi của thuộc tính theo thời gian (Slowly Changing Dimension).
\item Xác định bản ghi đại diện cho giá trị không xác định.
\end{enumerate}
\TaiLieu{Ch3 tr.22, 31--34. Khóa thay thế và Slowly Changing Dimension không có trong slide, xem Kimball \cite{kimball}.}
\KetQua{Bảng mô tả cột của từng bảng chiều.}
\PhanNhom{mô tả các bảng chiều}

\subsubsection{Dữ liệu dẫn xuất và bảng tổng hợp}
\CacBuoc
\begin{enumerate}
\item Liệt kê các trường dẫn xuất cần tính, dựa trên Bảng~\ref{tab:kpi}.
\item Xác định các truy vấn thường dùng cần tổng hợp sẵn.
\item Thiết kế các bảng tổng hợp hoặc data mart nếu cần.
\end{enumerate}
\TaiLieu{Ch2-cont tr.9, 13; Ch3 tr.7.}
\KetQua{Danh sách trường dẫn xuất và bảng tổng hợp.}
\PhanNhom{trường dẫn xuất và bảng tổng hợp}

\subsubsection{Sơ đồ lược đồ hoàn chỉnh}
\CacBuoc
\begin{enumerate}
\item Vẽ sơ đồ gồm toàn bộ bảng sự kiện, bảng chiều, khóa và quan hệ.
\item Kiểm tra lại sơ đồ có tính được toàn bộ chỉ số trong Bảng~\ref{tab:kpi} hay không.
\end{enumerate}
\TaiLieu{Ch3 tr.32, 34, 36 (các sơ đồ mẫu).}
\KetQua{Sơ đồ lược đồ hoàn chỉnh của kho dữ liệu.}
\HinhCho{Sơ đồ lược đồ của kho dữ liệu}{fig:schema}
